% \input{\pathsections "sec model"}

%%  %%  %%  %%  %%
\section{Model}
%
\begin{frame}\frametitle{The Model}
	%
\center
The \bl{model} is an \bl{input} to the least squares problem\\ along with the data.
	%
\end{frame}
%
\begin{frame}\frametitle{The Model}
	%
Let's \bl{derive the model} from first principles.
	%
\end{frame}
%

%
%%  %%  %%  %%  %%
\subsection{Heat Flow in 1 Dimension}
%
\begin{frame}\frametitle{\href{https://mathworld.wolfram.com/HeatConductionEquation.html}{Heat Conduction Equation}}
	%
Partial Differential Equation \\
\begin{equation}
	T_{t}(x,t) = \alpha \Delta T(x,t)
\label{eq:pde heat}
\end{equation}
\ \\[10pt]
Dirichlet Boundary Conditions\\
\begin{equation}
\begin{split}
	T(x=0,t) &= 0, \\
	T(x=100,t) &= 100,\\
%\label{eq:}
\end{split}
\end{equation}
	%
\end{frame}
%
%
\begin{frame}\frametitle{Stationary Solution for \eqnHeat}
\begin{equation}
	\lim\limits_{t\to\infty} T_{t}(x,t) = 0
\end{equation}
\end{frame}
%
%
\begin{frame}\frametitle{\eqnLaplace}
	%
The stationary solution for the \bl{\eqnHeat}...
\begin{equation}
	T_{t}(x,t) = \alpha \Delta T(x,t)
\label{eq:pde heat}
\end{equation}
	%
\begin{equation*}
	T_{t} = 0,
%\label{eq:}
\end{equation*}
\ \\
	%
	\pause
	%
...is given by \bl{\eqnLaplace}
\begin{equation}
	0 = \alpha \Delta T(x,t)
\label{eq:pde heat}
\end{equation}
	%
\end{frame}
%
%%  %%  %%  %%  %%
\subsection{\eqnLaplace}
%
\begin{frame}\frametitle{Solve Using Integration}
	%
\begin{equation}
	0 = \alpha \Delta T(x,t) , \qquad 0 \le x \le 10
\label{eq:pde heat}
\end{equation}
	%
\begin{equation}
	T(x) = 0 + 100 x
\label{eq:exact}
\end{equation}
%
\end{frame}
%
%
\begin{frame}\frametitle{Linear Model with Free Variables}
	%
Let's generalize the solution
\begin{equation}
\begin{split}
		%
	T(x) &= \bzero + \bl{100} x\\
		& \Downarrow \\
	T(x) &= \bl{a_{0}} + \bl{a_{1}} x
		%
\end{split}
%\label{eq:}
\end{equation}
\end{frame}


%%  %%  %%  %%  %%
\subsection{Linear Model}

\begin{frame}\frametitle{Model of Heat Flow}
\center
Temperature as a \bl{linear function} of position
	%
\begin{equation}
		%
	T(x) = \bazero + \baone x
		%
\label{eq:model}
\end{equation}
%
Solution parameters: \\
$\bazero$ intercept \\
$\baone$ slope
	%
\end{frame}

\begin{frame}\frametitle{Least Squares Inputs and Outputs}
	%
\begin{table}[htbp]
%\caption{Best fit to a linear function.}
	\begin{center}
		\begin{tabular}{ccc}
			%
			input & $\Rightarrow$ & output \\\hline
			%
			\includegraphics[ width = 2.0in ]{\pLocalGraphics bev-have} & \raisebox{1.8 cm}{$\Rightarrow$} &
			%
			\includegraphics[ width = 2.0in ]{\pLocalGraphics bev-want} \\
			%
			\{position, temperature\}$_{k=1}^{m}$ & $\Rightarrow$ & intercept and slope \\
			%
			$\lst{x_{k},T_{k}}_{k=1}^{m}$ 
			& $\Rightarrow$ & $\bl{\mat{c}{ a_{0} \\ a_{1} }} \pm \bl{\mat{c}{ \sigma_{0} \\ \sigma_{1}} }$\\
			%
		\end{tabular}
	\end{center}
\label{tab:bevington input-output}
\end{table}%
	%
\end{frame}

\begin{frame}\frametitle{Least Squares Inputs and Outputs}
	%
Inputs: \\
%
\begin{enumerate}
	\item Model function, $T(x) = \bazero + \baone x$
	\item Sequence of measurements, $\seqinputs{x}{T}$
\end{enumerate}
	\ \\[10pt]
	%
	\pause
	%
Outputs: Model parameters\\
%
\begin{enumerate}
	\item $\bazero \pm \bl{\sigma_{0}}$, intercept with error
	\item $\baone \pm \bl{\sigma_{1}}$, slope with error
\end{enumerate}
	%
\end{frame}


\endinput  %  ==  ==  ==  ==  ==  ==  ==  ==  ==
